% GENERATED FILE. Edit canonical lessons, then run npm run generate:books.
% canonical-source-hash: fnv1a64:8d843234a9468ef5
% canonical-lessons: ZH-C134-zhong, ZH-W134-zhong-heavy, ZH-C134-ruan, ZH-C134-ying, ZH-C134-kong, ZH-C134-man, ZH-R134-review-129-134

\chapter{Heavy and Soft}
\label{ch:zh-heavy-and-soft}
\hlchaptermodality{134}

\begin{chapteropening}
\textbf{By the end of this chapter:} \emph{I can say a thing is heavy, soft, hard, empty or full.}

Complete the last lesson of chapter 134: I can say a thing is heavy, soft, hard, empty or full.
\end{chapteropening}

\section[\texorpdfstring{\zh{重} (zhòng)}{重 (zhòng)}]{\zh{重} (zhòng) — heavy}

\label{lesson:ZH-C134-zhong}



\begin{quote}
Before the new one: say the Mandarin for to wipe, then the Mandarin for to fold.
\end{quote}

\subsection*{You'll want to know: \zh{重}}
\textbf{\zh{重}} — \emph{zhòng} — \textquotedblleft{}heavy\textquotedblright{}.

\begin{quote}
\textbf{\zh{这个太重了}} \emph{zhège tài zhòng le} — This one is too heavy.
\end{quote}

The new character, in pieces: \zh{重} — a short fall and a long bar on top, and under them \zh{里}, its vertical running up through the bar. The next lesson puts it in your hand.

\begin{cousinweb}
\textbf{\zh{重}} (\emph{zhòng}) is inherited from Middle Chinese \emph{drjowngX} (Guangyun, in Baxter's transcription), glossed 'much', 'thick' and 'careful'. 'Heavy' is not among the glosses, so the step to it is not settled.
\end{cousinweb}

\begin{grammarlens}[title={What you've built}]
One more word for describing things.
\end{grammarlens}

\subsection*{Your turn}
\begin{itemize}
\raggedright
  \item \emph{Say it:} \emph{zhòng}
  \item \emph{Say it:} \emph{zhòng}, once more
  \item \emph{Say it:} \emph{zhège tài zhòng le} — This one is too heavy.
  \item \emph{Recall:} say the Mandarin for to wipe, then the Mandarin for to fold, then say \emph{zhòng} again
\end{itemize}

\begin{tcolorbox}[breakable,colback=teal!4,colframe=teal!35!black,title={Before you move on}]
What does \zh{重} mean? (\textquotedblleft{}Heavy\textquotedblright{}.) Say it once more.
\end{tcolorbox}
\section[\texorpdfstring{\zh{重} (zhòng)}{重 (zhòng)}]{\zh{重} — a fall and a bar over \zh{里}}

\label{lesson:ZH-W134-zhong-heavy}



\begin{quote}
The word first, then the character inside it.

\begin{itemize}
\raggedright
  \item \emph{Recall:} say \emph{zhòng} — \textbf{R1}, one lesson back, before it decays
\end{itemize}

\begin{quote}
\zh{里} \zh{千}
\end{quote}

\emph{Read it:} each one aloud, with nothing said first — \emph{lǐ}, \emph{qiān}
\end{quote}

\subsection*{Script — a fall and a bar over \zh{里}}
\hlblockfigure{\detokenize{figures/ZH-W134-zhong-heavy-filmstrip.pdf}}{How \zh{重} is written, stroke by stroke}

\begin{quote}
\zh{重}
\end{quote}

\textbf{Two pieces, top to bottom.}

\textbf{A short left-falling stroke}, from the right down to the left, then \textbf{a long horizontal} across the whole width, as \zh{千} begins.

\textbf{\zh{里}, seven strokes}, in the order your hand knows: the box (its left side, the top that turns down the right side, the middle bar, the bar that closes it), then the vertical, which here starts above the long bar and runs down through it and the box, then the two bars of the foot, the lower one longer.

2 + 7 = nine.

\subsection*{Your turn}
\begin{itemize}
\raggedright
  \item \emph{Write it:} \zh{重}, all nine, with the model in front of you
  \item \emph{Cover:} write it again from memory, and count the pieces, not the strokes
\end{itemize}

\begin{tcolorbox}[breakable,colback=teal!4,colframe=teal!35!black,title={Before you move on}]
What sits under the long bar of \zh{重}? (\textbf{\zh{里}}.) Where does the vertical start? (\textbf{Above the long bar}.) Which is the last stroke? (\textbf{The long bar at the foot}.)

Source: \href{https://raw.githubusercontent.com/chanind/hanzi-writer-data/68d10a4b21150cae5e1ebbd223eed289cf32d90c/data/\%E9\%87\%8D.json}{Hanzi Writer Data \zh{重}, pinned stroke paths and medians}.
\end{tcolorbox}
\section[ruǎn]{ruǎn — soft}

\label{lesson:ZH-C134-ruan}



\begin{quote}
Before the new one: say the Mandarin for to fold, then the Mandarin for heavy.
\end{quote}

\subsection*{You'll want to know: ruǎn}
\emph{ruǎn} — \textquotedblleft{}soft\textquotedblright{}.

\begin{quote}
\emph{zhè zhāng chuáng hěn ruǎn} — This bed is soft.
\end{quote}

Its characters come later; for now it is a word for the ear.

\begin{cousinweb}
\emph{ruǎn} is inherited from Middle Chinese \emph{nywenX} (Guangyun, in Baxter's transcription), glossed 'soft', under an older way of writing the character.
\end{cousinweb}

\begin{grammarlens}[title={What you've built}]
One more word for describing things.
\end{grammarlens}

\subsection*{Your turn}
\begin{itemize}
\raggedright
  \item \emph{Say it:} \emph{ruǎn}
  \item \emph{Say it:} \emph{ruǎn}, once more
  \item \emph{Say it:} \emph{zhè zhāng chuáng hěn ruǎn} — This bed is soft.
  \item \emph{Recall:} say the Mandarin for to fold, then the Mandarin for heavy, then say \emph{ruǎn} again
\end{itemize}

\begin{tcolorbox}[breakable,colback=teal!4,colframe=teal!35!black,title={Before you move on}]
What does ruǎn mean? (\textquotedblleft{}Soft\textquotedblright{}.) Say it once more.
\end{tcolorbox}
\section[yìng]{yìng — hard; stiff}

\label{lesson:ZH-C134-ying}



\begin{quote}
Before the new one: say the Mandarin for heavy, then the Mandarin for soft.
\end{quote}

\subsection*{You'll want to know: yìng}
\emph{yìng} — \textquotedblleft{}hard; stiff\textquotedblright{}.

\begin{quote}
\emph{zhège miànbāo tài yìng le} — This bread is too hard.
\end{quote}

Its characters come later; for now it is a word for the ear.

\begin{cousinweb}
\emph{yìng} is inherited from Middle Chinese \emph{ngeangH} (Guangyun, in Baxter's transcription), glossed 'firm and sturdy'.
\end{cousinweb}

\begin{grammarlens}[title={What you've built}]
One more word for describing things.
\end{grammarlens}

\subsection*{Your turn}
\begin{itemize}
\raggedright
  \item \emph{Say it:} \emph{yìng}
  \item \emph{Say it:} \emph{yìng}, once more
  \item \emph{Say it:} \emph{zhège miànbāo tài yìng le} — This bread is too hard.
  \item \emph{Recall:} say the Mandarin for heavy, then the Mandarin for soft, then say \emph{yìng} again
\end{itemize}

\begin{tcolorbox}[breakable,colback=teal!4,colframe=teal!35!black,title={Before you move on}]
What does yìng mean? (\textquotedblleft{}Hard; stiff\textquotedblright{}.) Say it once more.
\end{tcolorbox}
\section[kōng]{kōng — empty}

\label{lesson:ZH-C134-kong}



\begin{quote}
Before the new one: say the Mandarin for soft, then the Mandarin for hard to the touch.
\end{quote}

\subsection*{You'll want to know: kōng}
\emph{kōng} — \textquotedblleft{}empty\textquotedblright{}.

\begin{quote}
\emph{bīngxiāng shì kōng de} — The fridge is empty.
\end{quote}

Its characters come later; for now it is a word for the ear.

\begin{cousinweb}
\emph{kōng} is inherited from Middle Chinese \emph{khuwng} (Guangyun, in Baxter's transcription), glossed 'empty, hollow'.
\end{cousinweb}

\begin{grammarlens}[title={What you've built}]
One more word for describing things.
\end{grammarlens}

\subsection*{Your turn}
\begin{itemize}
\raggedright
  \item \emph{Say it:} \emph{kōng}
  \item \emph{Say it:} \emph{kōng}, once more
  \item \emph{Say it:} \emph{bīngxiāng shì kōng de} — The fridge is empty.
  \item \emph{Recall:} say the Mandarin for soft, then the Mandarin for hard to the touch, then say \emph{kōng} again
\end{itemize}

\begin{tcolorbox}[breakable,colback=teal!4,colframe=teal!35!black,title={Before you move on}]
What does kōng mean? (\textquotedblleft{}Empty\textquotedblright{}.) Say it once more.
\end{tcolorbox}
\section[mǎn]{mǎn — full}

\label{lesson:ZH-C134-man}



\begin{quote}
Before the new one: say the Mandarin for hard to the touch, then the Mandarin for empty.
\end{quote}

\subsection*{You'll want to know: mǎn}
\emph{mǎn} — \textquotedblleft{}full\textquotedblright{}.

\begin{quote}
\emph{wǒ de bēizi mǎn le} — My cup is full.
\end{quote}

Its characters come later; for now it is a word for the ear.

\begin{cousinweb}
\emph{mǎn} is inherited from Middle Chinese \emph{manX} (Guangyun, in Baxter's transcription), glossed 'full' and 'to fill'.
\end{cousinweb}

\begin{grammarlens}[title={What you've built}]
One more word for describing things.
\end{grammarlens}

\subsection*{Your turn}
\begin{itemize}
\raggedright
  \item \emph{Say it:} \emph{mǎn}
  \item \emph{Say it:} \emph{mǎn}, once more
  \item \emph{Say it:} \emph{wǒ de bēizi mǎn le} — My cup is full.
  \item \emph{Recall:} say the Mandarin for hard to the touch, then the Mandarin for empty, then say \emph{mǎn} again
\end{itemize}

\begin{tcolorbox}[breakable,colback=teal!4,colframe=teal!35!black,title={Before you move on}]
What does mǎn mean? (\textquotedblleft{}Full\textquotedblright{}.) Say it once more.
\end{tcolorbox}
\section[(dialogue)]{Looking back: from Quiet and Noisy to Heavy and Soft}

\label{lesson:ZH-R134-review-129-134}



\begin{quote}
Say the Mandarin for quiet and noisy.
\end{quote}

\subsection*{Your turn}
\begin{itemize}
\raggedright
  \item \emph{Say it:} quiet, noisy, bustling, convenient, interesting
  \item \emph{Say it:} a sick person, to catch a cold, to have a fever, sore, to see a doctor
  \item \emph{Say it:} spring, summer, autumn, winter, season
  \item \emph{Say it:} thousand, ten thousand, zero, in all, to count
  \item \emph{Say it:} to clean (a room), housework, to tidy up, to wipe, to fold
  \item \emph{Say it:} heavy, soft, hard to the touch, empty, full
  \item \emph{Read it:} \zh{吵} \zh{病} \zh{春} \zh{千} \zh{扫} \zh{重} — each one aloud, then a word that holds it
\end{itemize}

\begin{tcolorbox}[breakable,colback=teal!4,colframe=teal!35!black,title={Before you move on}]
Quiet? (\emph{ānjìng}.) Noisy? (\textbf{\zh{吵}}, \emph{chǎo}.) Bustling? (\emph{rènao}.) Convenient? (\emph{fāngbiàn}.) Interesting? (\emph{yǒuyìsi}.) A sick person? (\textbf{\zh{病人}}, \emph{bìngrén}.) To catch a cold? (\emph{gǎnmào}.) To have a fever? (\emph{fāshāo}.) Sore? (\emph{téng}.) To see a doctor? (\textbf{\zh{看病}}, \emph{kànbìng}.) Spring? (\textbf{\zh{春天}}, \emph{chūntiān}.) Summer? (\emph{xiàtiān}.) Autumn? (\emph{qiūtiān}.) Winter? (\emph{dōngtiān}.) Season? (\emph{jìjié}.) Thousand? (\textbf{\zh{千}}, \emph{qiān}.) Ten thousand? (\emph{wàn}.) Zero? (\emph{líng}.) In all? (\emph{yígòng}.) To count? (\emph{shǔ}.) To clean (a room)? (\textbf{\zh{打扫}}, \emph{dǎsǎo}.) Housework? (\emph{jiāwù}.) To tidy up? (\emph{shōushi}.) To wipe? (\emph{cā}.) To fold? (\emph{dié}.) Heavy? (\textbf{\zh{重}}, \emph{zhòng}.) Soft? (\emph{ruǎn}.) Hard to the touch? (\emph{yìng}.) Empty? (\emph{kōng}.) Full? (\emph{mǎn}.)
\end{tcolorbox}
